\subsection{The natural first torsion group of the free-disk partial model}

\paragraph{A-S09-18 (Proved; a natural description of degree-four torsion).}
Let $P$ be a strict PD cochain $S$-algebra, where
$S=\Gamma_{\mathbb Z}(T)$, which is free over $S$
on any collection of vertical identity disks
\[
                  dx_a=y_a,
 \qquad |x_a|=q_a\geq0,\quad |y_a|=q_a+1,
 \qquad \operatorname{wt}(x_a)=\operatorname{wt}(y_a)>0.
 \tag{S09.78}
\]
All graded commutativity is alternating, and the even
generators have the specified divided powers.
Let $P^{>0,0}$ denote its positive-weight zero-forms,
and define the weight-graded abelian group
\[
 Q^1(P)=P^1\big/
             \bigl(dP^0+P^{>0,0}P^1\bigr).
 \tag{S09.79}
\]
There is a canonical isomorphism
\[
 H^4(P)^w\cong
 \bigoplus_{s>0}
 (S^{w-2s}/2S^{w-2s})\otimes_{\mathbb F_2}
                              (Q^1(P)^s/2Q^1(P)^s),
 \qquad
 a\otimes[x]\longmapsto a[\gamma_2(dx)].             \tag{S09.80}
\]
Terms with $w-2s<0$ are zero.  This isomorphism is
natural for every coefficient-preserving strict PD
cochain map between algebras of the stated kind;
the maps need not preserve any chosen disk generators
or chosen decomposition.  In particular it is an
isomorphism of simplicial weight-graded groups when
each simplex algebra has such a free-disk decomposition.

\begin{proof}
For an odd one-form $x$, its differential is an even
two-form, so $\gamma_2(dx)$ is a closed four-form.
The identities
\[
 \begin{split}
 \gamma_2(d(x+x'))-\gamma_2(dx)-\gamma_2(dx')
                          &=d(x\,dx'),\\
 2\gamma_2(dx)&=d(x\,dx),\\
 \gamma_2(d(da))&=0
 \end{split}                                                   \tag{S09.81}
\]
show additivity in cohomology, annihilation by two,
and invariance under adding an exact one-form.
For a positive-weight zero-form $a$, which belongs
to the PD ideal but need not be closed, the exact
nonclosed-factor identity is
\[
              \gamma_2(d(a x))
                         =d\bigl(\gamma_2(a)x\,dx\bigr).
 \tag{S09.82}
\]
This is the $n=2$ instance of A-S09-13 and can also
be checked directly: the two terms of the derivative
are $a\,da\,x\,dx$ and $a^2\gamma_2(dx)$, precisely
the divided-power expansion on the left.  Therefore
(S09.81)--(S09.82) show that the formula on the
right of (S09.80) factors through (S09.79) modulo
two and is independent of every representative.
They also prove its naturality under every stated
strict PD cochain map.  Weight is doubled on the
parameter because $\gamma_2$ has arity two.

It remains to prove that the map is an isomorphism.
Choose a free-disk presentation only for this check.
In degree one, the quotient (S09.79) kills all
positive-weight zero-form coefficients.  It also
kills the odd generators $y_a=dx_a$ with $q_a=0$.
Thus $Q^1(P)$ is freely generated as an abelian
group by exactly the $x_a$ with $q_a=1$.

An even-to-odd disk is the direct sum of its constants
and identity pairs
$\gamma_m(x)\mapsto\gamma_{m-1}(x)y$.
An odd-to-even disk is the direct sum of its constants
and the pairs
\[
 x\gamma_{m-1}(y)\xrightarrow{\ m\ }\gamma_m(y),
          \qquad\text{in degrees }m(q+1)-1,m(q+1).
 \tag{S09.83}
\]
In a finite tensor product, every summand containing
an identity pair contracts over the remaining factors.
The first possible nonunit pair is $q=1,m=2$, in
degrees three and four, with cokernel $\mathbb Z/2$.
Two nonunit pairs have minimum total degree six,
and every other single nonunit pair has its cokernel
in degree at least six.  Hence the full degree-four
cohomology is a direct sum of one copy of $S/2S$
for each disk with $q_a=1$, represented precisely
by $\gamma_2(y_a)$.  This agrees with the stated
basis of (S09.79) and proves (S09.80).  Infinite
collections follow from exact filtered colimits
of finite tensor products; each element uses only
finitely many generators.  No basis choice remains
in the already defined natural map.
\end{proof}

\paragraph{A-S09-19 (Proved; the C21 partial source can retain integral exactness in simplices zero and one).}
For every pair of targets to which C21 applies, its
common partial source and its two actual strict maps
can be chosen with $P_0=S$ and with
$S\to P_1$ an integral quasi-isomorphism in every
weight.  All the conclusions and uniform bounds of
C21 remain unchanged: terms are free, the horizontal
bound is $\mu_P(d)=\max_{0\leq q\leq d}W(q)$ for
the same permitted choices of $W(q)$, the stated
factorial-local and prime-local vertical ranges
hold, and $q\leq w$ is retained when imposed.

\begin{proof}
In the actual C21 construction, omit zero normalized
cycles when attaching horizontal cells.  A cell
$K(n)\hookrightarrow D(n+1)$ has quotient $K(n+1)$,
so it never introduces a simplex-zero generator.
Consequently $P_0=S$.  Every normalized degree-zero
cycle in a positive form row is then zero, so no
positive-form cell with $n=0$ is needed.  All such
cells have quotient parameter zero in simplex one.
Thus $P_1$ is free over $S$ only on form-zero
even-to-odd disks, and its coefficient inclusion
is an integral quasi-isomorphism.  Their contractions
are the integral identity pairs in the proof above,
and tensoring them gives the asserted coefficient
comparison, including every weight.

Omitting zero cycles loses no killing step: zero
already has the zero primitive.  Every nonzero
normalized cycle in a round is still attached at
the next round, and every cycle in the filtered
colimit is represented in some finite round.
Consequently all of the original horizontal
exactness conclusions hold with exactly the same
stage weights $W(q)$.  The two target maps extend
through each retained cell by the identical
chosen target primitives used in C21.  Every
simplex algebra remains a free extension by
vertical identity disks, so the original freeness,
support, and vertical disk decomposition proofs
apply without change.  This verifies every stated
conclusion by the actual retained cell maps.
\end{proof}

Formula (S09.79)
gives, more specifically,
\[
                  Q^1(P_0)=Q^1(P_1)=0,
             \qquad H^4(P)_0=H^4(P)_1=0.            \tag{S09.84}
\]
This is a pointwise vanishing statement in low
simplicial degrees, including every weight.

The formula is deliberately stated for the exact
free identity-disk simplex algebras to which its
proof applies.  A later two-cell repair can introduce
nonunit disks or differentials with old-algebra
terms; no identification of its new cohomology
with (S09.80) is inferred.  The usefulness of the
formula for the initial source is that every
nonlinear simplicial substitution has already been
checked by (S09.81)--(S09.82), so the coefficient
and parameter comparison in (S09.80) is an actual
natural map rather than a choice of linear coordinates.
